\textbf{Lombard \& Piraux 2004 --- Eq 32--37 (ESIM Taylor setup, paper pp.\ 13--14)}

\medskip

\textbf{Setup (Fig 3).} Let $M(x_I, y_J) \in \Omega_2$ be an
\emph{irregular} grid point (whose stencil reaches across $\Gamma$),
and $P$ its orthogonal projection onto $\Gamma$. Let $B$ be the set
of grid points enclosed in the disk of radius $q$ centred on $P$,
partitioned $B = B_1 \cup B_2$ by side.

\medskip

\textbf{Eq 32 --- $k$-th order Taylor row vector $\boldsymbol{\Pi}_{i,j}^k$:}
\begin{equation*}
  \boldsymbol{\Pi}_{i,j}^k =
  \left( 1,\ (x_i - x_P),\ (y_j - y_P),\ \ldots,\ \frac{(y_j - y_P)^k}{k!} \right)
  \tag{32}
\end{equation*}

\textbf{Eq 33 --- modified value at $M(x_I, y_J)$:}
\begin{equation*}
  \boldsymbol{U}^*(x_I, y_J, t^n)
  = \boldsymbol{\Pi}_{I,J}^k \boldsymbol{U}_1^k
  \tag{33}
\end{equation*}

\medskip

\textbf{Eq 34 --- Taylor expansion at $(i, j) \in B_1$} (using Eq 25, 26):
\begin{equation*}
\begin{aligned}
  (i,j) \in B_1,\ \boldsymbol{U}(x_i, y_j, t^n)
  &= \boldsymbol{\Pi}_{i,j}^k \boldsymbol{U}_1^k + O(\Delta x^{k+1}) \\
  &= \boldsymbol{\Pi}_{i,j}^k \boldsymbol{G}_1^k \boldsymbol{V}_1^k + O(\Delta x^{k+1}) \\
  &= \boldsymbol{\Pi}_{i,j}^k \boldsymbol{G}_1^k \boldsymbol{K}_1^k \boldsymbol{W}_1^k + O(\Delta x^{k+1})
\end{aligned}
\tag{34}
\end{equation*}

\textbf{Eq 35 --- inject the $[\boldsymbol{I} | \boldsymbol{0}]$ projector
to use the full $(\boldsymbol{W}_1^k, \boldsymbol{\Lambda}^k)$ vector:}
\begin{equation*}
  (i,j) \in B_1,\ \boldsymbol{U}(x_i, y_j, t^n)
  = \boldsymbol{\Pi}_{i,j}^k \boldsymbol{G}_1^k \boldsymbol{K}_1^k
  (\boldsymbol{I} \mid \boldsymbol{0})
  \begin{pmatrix} \boldsymbol{W}_1^k \\ \boldsymbol{\Lambda}^k \end{pmatrix}
  + O(\Delta x^{k+1})
  \tag{35}
\end{equation*}

\medskip

\textbf{Eq 36 --- Taylor expansion at $(i, j) \in B_2$:}
\begin{equation*}
\begin{aligned}
  (i,j) \in B_2,\ \boldsymbol{U}(x_i, y_j, t^n)
  &= \boldsymbol{\Pi}_{i,j}^k \boldsymbol{U}_2^k + O(\Delta x^{k+1}) \\
  &= \boldsymbol{\Pi}_{i,j}^k \boldsymbol{G}_2^k \boldsymbol{V}_2^k + O(\Delta x^{k+1}) \\
  &= \boldsymbol{\Pi}_{i,j}^k \boldsymbol{G}_2^k \boldsymbol{K}_2^k \boldsymbol{W}_2^k + O(\Delta x^{k+1})
\end{aligned}
\tag{36}
\end{equation*}

\textbf{Eq 37 --- substitute Eq 31 (express $\boldsymbol{W}_2^k$
via $\boldsymbol{W}_1^k$ and $\boldsymbol{\Lambda}^k$):}
\begin{equation*}
\begin{aligned}
  (i,j) \in B_2,\ \boldsymbol{U}(x_i, y_j, t^n)
  &= \boldsymbol{\Pi}_{i,j}^k \boldsymbol{G}_2^k \boldsymbol{K}_2^k
    \boldsymbol{S}_2^{k\,-1}
    (\boldsymbol{S}_1^k \mid \boldsymbol{R}_{S_2}^k)
    \begin{pmatrix} \boldsymbol{W}_1^k \\ \boldsymbol{\Lambda}^k \end{pmatrix} \\
  &\phantom{=}\ + O(\Delta x^{k+1})
\end{aligned}
\tag{37}
\end{equation*}

\medskip

\textbf{Note:} Eq 32--33 are on paper p.\ 13 (bottom);
Eq 34--37 are on paper p.\ 14. The combined paper region shown on
the left of this PDF spans both pages.

\textbf{Next.} Eq 38--42 (paper p.\ 15) assemble these row-by-row
Taylor expansions into the over-determined matrix $\boldsymbol{M}$,
compute its least-squares pseudoinverse, and arrive at the master
modified-value equation
$\boldsymbol{U}^*_{I,J} = \boldsymbol{\Pi}_{I,J}^k \boldsymbol{G}_1^k
\boldsymbol{K}_1^k \bar{\boldsymbol{M}}^{-1} (\boldsymbol{U}^n)_B$.
